\title{Casper the Friendly Finality Gadget}
\author{
        Vitalik Buterin and Virgil Griffith \\
        Ethereum Foundation}

\documentclass[12pt, final]{article}

%% make in NIPS format
\usepackage{nips10submit_e}
\nipsfinalcopy



\usepackage{graphicx}
\DeclareGraphicsExtensions{.pdf,.eps,.png,.jpg}		% search for .pdf, then .eps, then .pngs, then .jpg

% look in these subdirectories for graphics referenced by \includegraphics
% each entry must end with a /
\graphicspath{{figs/}{figures/}{images/}{./}}

\newcommand*{\red}[1]{ \color{red} #1}

\usepackage{tabularx}
\usepackage{url}
\usepackage{cite}
\usepackage{amsmath}

% for some math symbols.  Particularly \centerdot
\usepackage{ amssymb }


%\usepackage{ lmodern }			% I prefer this over times

\usepackage{array}			% replacement for eqnarray.  Must be BEFORE \usepackage{arydshln}
\usepackage{units}			% for \nicefrac{\alpha}{\beta}


\usepackage{amsthm}		% for theorems
\newtheorem{definition}{Definition}

%\usepackage{microtype}

%\usepackage{wasysym}

\usepackage{textcomp, marvosym} % pretty symbols
\usepackage{booktabs} 	% for much better looking tables

% for indicator functions
\usepackage{dsfont}

% For automatic capitalizaton of section names, etc.
\usepackage{titlesec,titlecaps}


\Addlcwords{is with of the and in}
%\Addlcwords{of the}
%\Addlcwords{and}
\titleformat{\section}[block]{}{\normalfont\Large\bfseries\thesection.\;}{0pt}{\formatsectiontitle}
\newcommand{\formatsectiontitle}[1]{\normalfont\Large\bfseries\titlecap{#1}}

\titleformat{\subsection}[block]{}{\normalfont\large\bfseries\thesubsection.\;}{0pt}{\formatsubsectiontitle}
\newcommand{\formatsubsectiontitle}[1]{\normalfont\large\bfseries\titlecap{#1}}





% for pretty Euler script
\usepackage[mathscr]{euscript}
\usepackage{bold-extra}



\usepackage{fullpage} % save trees

\usepackage{subfig}
%\usepackage{float} % for \subfloat

%%%%%%%%%%%%%%%%%%%%%%%%%%%%%%%%%%%%%%%%%%
%% More customizable Lists
%%%%%%%%%%%%%%%%%%%%%%%%%%%%%%%%%%%%%%%%%%
% Better symbols custom enumerative lists, define any symbol you'd like
\usepackage{enumitem}


%%%%%%%%%%%%%%%%%%%%%%%%%%%%%%%%%%%%%%%%%%
%% Custom Symbols 
%%%%%%%%%%%%%%%%%%%%%%%%%%%%%%%%%%%%%%%%%%
% \xspace at the end of custom macros never fucks up spacing.
% example of best practice: \newcommand{\apples}{\textsf{AppLeS}\xspace} 
\usepackage{xspace}



%%%%%%%%%%%%%%%%%%%%%%%%%%%%%%%%%%%%%%%%
%% Abbreviations you'll always want
%%%%%%%%%%%%%%%%%%%%%%%%%%%%%%%%%%%%%%%%
\newcommand*{\TODO}[1]{{\centering {\sffamily \color{red} #1} \vskip10pt }}
\newcommand*{\todo}[1]{{\sffamily [{\color{red} #1}]}}
\newcommand*{\fix}[1]{{\sffamily [{\textnormal{\color{red} #1}}]}}



%-----------------------------------------------------------------------------
%  Cross references
%-----------------------------------------------------------------------------
% The following code defines how you make references to figures, tables, etc...
% It is defined in one place only, and can be modified for all references
% in the document at the same time.
% Instead of typing each time: "see Fig. \ref{myfig}" you can create a command
% \figref which will do the job. Then in text you only type \figref{myfig} and LaTeX
% will do the rest.
\newcommand{\tblref}[1]{Table~\ref{#1}}
%\renewcommand*{\figref}[1]{Fig.~\ref{#1}}
\renewcommand{\eqref}[1]{eq.~\ref{#1}}


\newcommand{\figref}[1]{Figure~\ref{#1}}
\newcommand{\Figref}[1]{Figure~\ref{#1}}
\newtheorem{theorem}{Theorem}
\newtheorem{lemma}[theorem]{Lemma}

%%%%%%%%%%%%%%%%%%%%%%%%%%%%%%%%%%%%%%%%

% for \sout{} for strikeout
\usepackage[normalem]{ulem}


% for better manipulation of tables
\usepackage{makecell}
\renewcommand\theadfont{\bfseries}


%-----------------------------------------------------------------------------
%  Misc symbols that I like
%-----------------------------------------------------------------------------
\newcommand*{\opname}[1]{\operatorname{#1}}


\renewcommand*{\to}{\rightarrow}


%%%%





\usepackage{color}
\newcommand*{\todo}[1]{\color{red} #1}
%\newcommand{\eqref}[1]{eq.~\ref{#1}}
\newcommand{\figref}[1]{Figure~\ref{#1}}
\usepackage{amsthm}
\usepackage{amsmath}
\newtheorem{theorem}{Theorem}
\newtheorem{definition}{Definition}
\usepackage{graphicx}
\graphicspath{{figs/}{figures/}{images/}{./}}

%% Special symbols we'll probably iterate on
%%%%%%%%%%%%%%%%%%%%%%%%%%%%%%%%%%%%%%%%%%%%%%%%%%%%%%
% we will probably iterate on these symbols until we have a notation we like
\newcommand{\epoch}{\ensuremath{e}\space}
\newcommand{\hash}{\textnormal{h}\space}

% symbols for the epoch and hash source
\newcommand{\epochsource}{\ensuremath{\epoch_{\star}}\space}
\newcommand{\hashsource}{\ensuremath{\hash_{\star}}\space}

\newcommand{\signature}{\ensuremath{\mathcal{S}}\space}

\newcommand{\totaldeposit}{\textnormal{TD}\space}

\newcommand{\gamesymbol}{\reflectbox{G}}

\newcommand{\msgPREPARE}{\textbf{\textsc{prepare}}\space}
\newcommand{\msgCOMMIT}{\textbf{\textsc{commit}}\space}

% Symbols for the Last Justified Epoch and Hash
\newcommand{\epochLJ}{\ensuremath{\epoch_{\textnormal{LJ}}}\space}
\newcommand{\hashLJ}{\ensuremath{\hash_{\textnormal{LJ}}}\space}

% Symbols for the Last Finalized Epoch and Hash
\newcommand{\epochLF}{\ensuremath{\epoch_{\textnormal{LF}}}\space}
\newcommand{\hashLF}{\ensuremath{\hash_{\textnormal{LF}}}\space}

% Griefing Factor symbol
\newcommand{\GF}[1]{\mathds{GF}\left( #1 \right)\space}

% Genesis block symbol
\newcommand{\Genesisblock}{\ensuremath{G}\space}

\begin{document}

\maketitle

\begin{abstract}
We give an introduction to the consensus algorithm details of Casper: the Friendly Finality Gadget, as an overlay on an existing proof of work blockchain such as Ethereum. Casper is a partial consensus mechanism inspired by a combination of existing proof of stake algorithm research and Byzantine fault tolerant consensus theory, which if overlaid onto another blockchain (which could theoretically be proof of work or proof of stake) adds strong \textit{finality} guarantees that improve the blockchain's resistance to transaction reversion (or ``double spend'') attacks.
\end{abstract}

%%%%%%%%%%%%%%%%%%%%%%%%%%%%%%%%%%%%%%%%%%%%%%%%%%%%%%%%%%%%%%%%%%%%%%%%%%%%%%%%%%%%%%%%%%%%%%%%%%%%%%%%%%%%%%%%%%%%%%%%%%
%%%%%%%%%%%%%%%%%%%%%%%%%%%%%%%%%%%%%%%%%%%%%%%%%%%%%%%%%%%%%%%%%%%%%%%%%%%%%%%%%%%%%%%%%%%%%%%%%%%%%%%%%%%%%%%%%%%%%%%%%%
\section{Introduction}
\label{sect:intro}

In the past few years there has been considerable research into ``proof of stake''-based blockchain consensus algorithms. In a proof of stake system, a blockchain grows and agrees on new blocks through a process where anyone who holds coins inside of the system can participate, and the amount of influence that any given coin holder has is proportional to the number of coins (or ``stake'') that they hold. This represents an alternative to proof of work ``mining'', allowing blockchains to operate without the high hardware and electricity costs that proof of work blockchains require.

There are two major schools of thought in proof of stake algorithm design. The first, \textit{chain-based proof of stake}, tries to closely mimic the mechanics of proof of work, featuring a chain of blocks and an algorithm that ``simulates'' mining by pseudorandomly assigning the right to create new blocks to stakeholders. This includes Peercoin\cite{king2012ppcoin}, Blackcoin\cite{vasin2014blackcoin} and Iddo Bentov's work\cite{bentov2016pos}.

The other school, \textit{BFT-based proof of stake}, is based on a thirty year old body of research into \textit{Byzantine fault tolerant consensus algorithms} such as PBFT \cite{castro1999practical}. BFT algorithms tend to have strong and rigorously proven mathematical properties; for example, one can usually mathematically prove that as long as at more than $\frac{2}{3}$ of participants in the protocol are following the protocol correctly, then the algorithm cannot possibly finalize conflicting block hashes at the same time (``safety''); this result holds regardless of network latency. The idea of repurposing BFT algorithms for proof of stake was first introduced by Tendermint\cite{kwon2014tendermint}.

\subsection{Our Work}

We follow the BFT tradition, though with some modifications. Casper the Friendly Finality Gadget is an \textit{overlay} atop a \textit{proposal mechanism}---a mechanism which proposes \textit{checkpoints} (this is similar to the common BFT abstraction of ``leader election'' except it is more complex so as to accommodate Casper's chain-based nature).  Casper is responsible for \textit{finalizing} these checkpoints.  Casper provides safety, but does not guarantee liveness by itself---Casper depends on the proposal mechanism for liveness.  That is, even if the proposal mechanism is wholly controlled by attackers, Casper prevents attackers from finalizing two conflicting checkpoints, however, the attackers can prevent Casper from finalizing any future checkpoints.

The proposal mechanism will initially be the existing Ethereum proof of work chain, making the first version of Casper a \textit{hybrid PoW/PoS algorithm} that relies on proof of work for liveness but not safety, but in future versions the proposal mechanism can be substituted with something else.

Our algorithm introduces several new properties that BFT algorithms do not necessarily support.
\begin{itemize}
\item We flip the emphasis of the proof statement from the traditional ``as long as $>\frac{2}{3}$ of validators are honest, there will be no safety failures'' to the contrapositive ``if there is a safety failure, then $\ge \frac{1}{3}$ of validators violated some protocol rule.''

\item We add \textit{accountability}.  If a validator violates the rules, we can detect the violation, and know who violated the rule: ``$\ge \frac{1}{3}$ violated the rules, \textit{and we know who they are}''.  Accountability allows us to penalize malfeasant validators, solving the \textit{nothing at stake} problem\cite{} that often plagues chain-based PoS. The penalty is the malfeasant validator's entire deposit.  This maximum penalty provides a bulwark against violating the protocol by making violations immensely expensive.  The economic robustness of protocol guarantees is much higher than the size of the rewards that the system pays out during normal operation.  This provides a \textit{much stronger} security guarantee than possible with proof of work.

\item We introduce a provably safe way for the validator set to change over time (Section \ref{sect:join_and_leave}).
\item We introduce a way to recover from attacks where more than $\frac{1}{3}$ of validators drop offline, at the cost of a very weak \textit{tradeoff synchronicity assumption} (Section \ref{sect:leak}).
\item The design of the algorithm as an overlay makes it easier to implement as an upgrade to an existing proof of work chain.
\end{itemize}

We will describe the protocol in stages, starting with a simple version (Section \ref{sect:protocol}) and then progressively adding features such as validator set changes (Section \ref{sect:join_and_leave}) and mass liveness fault recovery.

\section{The Casper Protocol}
\label{sect:protocol}
In the simple version, we assume there is a set of validators and a \textit{proposal mechanism} which is any system that proposes blocks (such as a proof of work chain). Under normal operation, almost all of the blocks proposed by the proposal mechanism would form a chain, but if the proposal mechanism operation is faulty (in proof of work, this may arise due to high latency or 51\% attacks) there may be multiple divergent chains being extended at the same time.


We order a chain of blockhashes into a sequence called a \emph{blockchain} $\mathbf{B}$.  Within blockchain $\mathbf{B}$ is there is a subset called \emph{checkpoints},

\begin{equation}
\begin{split}
    \mathbf{B} &\equiv \left( b_0, b_1, b_2, \ldots \right) \\
    \mathbf{C} &\equiv \left( b_0, b_{99}, b_{199}, b_{299}, \ldots \right) \;
\end{split}
\end{equation}

This leads to the formula for an arbitrary checkpoint,
\begin{equation}
    C_i = \begin{cases}
     b_0 \qquad \qquad \qquad\ \  \textnormal{if } i = 0, \\
     b_{ 100*i - 1 } \qquad \qquad \textnormal{otherwise.}
     \end{cases}
\end{equation}

$C_0$ is called the \textit{genesis}. An \emph{epoch} is defined as the contiguous sequence of blocks between two checkpoints, including the later checkpoint but not the earlier one.  The \textit{epoch of a block} is the index of the epoch containing that hash, e.g., the epoch of block 599 is 5. \footnote{To get the epoch of a particular block $b_i$, it is $epoch(b_i) = \lfloor i / 100 \rfloor$.}

Each validator has a \emph{deposit}; when a validator joins, their deposit is the number of coins that they deposited, and from there on each validator's deposit rises and falls with rewards and penalties. For the rest of this paper, when we say ``$\frac{2}{3}$ of validators'', we are referring to a \emph{deposit-weighted} fraction; that is, a set of validators whose sum deposit size equals to at least $\frac{2}{3}$ of the total deposit size of the entire set of validators. ``$\frac{2}{3}$ prepares'' will be used as shorthand for ``prepares from $\frac{2}{3}$ of validators''.

Validators can broadcast two types of messages:

$$\langle \msgPREPARE, \hash, \epoch, \hashsource, \epochsource, \signature \rangle$$


	\begin{tabular}{l l}
	\textbf{Notation} & \textbf{Description} \\
	\hash & a checkpoint hash \\
	\epoch & the epoch of the checkpoint \\
	$\hashsource$ & the most recent justified hash \\
	$\epochsource$ & the epoch of $\hashsource$  \\
	\signature & signature of $(\hash,\epoch,\hashsource,\epochsource)$ from the validator's private key \\
	\end{tabular} \label{tbl:prepare}

$$ $$

$$\langle \msgCOMMIT, \hash, \epoch, \signature \rangle$$
	
	\begin{tabular}{l l}
	\textbf{Notation} & \textbf{Description} \\
	\hash & a checkpoint hash \\
	\epoch & the epoch of the checkpoint \\
	\signature & signature from the validator's private key \\
	\end{tabular} 
	\label{tbl:commit}
	\label{fig:messages}

$$ $$

A checkpoint $\hash$ is considered \emph{justified} if there exists some \hashsource such that:

\begin{enumerate}
\item There exist $\frac{2}{3}$ prepares of the form $\langle \msgPREPARE, \hash, \epoch(\hash), \hashsource, \epoch(\hashsource), \signature \rangle$
\item \hashsource itself is justified
\end{enumerate}

Note that all $\frac{2}{3}$ prepares that justify \hash must have the same \hashsource; half with $\hashsource^1$ and half with $\hashsource^2$ will not suffice. A checkpoint \hash is considered \emph{finalized} if:

\begin{enumerate}
\item \hash is justified
\item There exist $\frac{2}{3}$ commits of the form $\langle \msgCOMMIT, \hash, \epoch(\hash), \signature \rangle$
\end{enumerate}

The genesis is considered to be justified and finalized, serving as the base case for the recursive definition.


During epoch $n$, validators are expected to send prepare and commit messages with $\epoch = n$ and $h$ equal to a checkpoint of epoch $n$. Prepare messages may specify as \hashsource a checkpoint for any previous epoch (preferably the preceding checkpoint) of \hash, and which is \textit{justified} (see below), and the \epochsource is expected to be the epoch of that checkpoint.

Validators only recognize prepares and commits that have been included in blocks (even if those blocks are not part of the main chain). \footnote{This simplifies our finalty mechanism because it allows it to be expressed as a fork choice rule where the ``score'' of a block only depends on the block and its descendants, similarly to the longest chain rule and GHOST\cite{sompolinsky2013accelerating} in proof of work. Also, because the blockchain will not accept (i) commits that do not point to an already justified checkpoint, (ii) prepares that do not point to an already justified \hashsource value, or (iii) prepares or commits that supply incorrect epoch numbers, this makes two of the four slashing conditions in earlier versions of Casper\cite{minslashing} superfluous.}

\subsection{Casper Commandments}

The most notable property of Casper is that it is impossible for two conflicting checkpoints to be finalized unless $\geq \frac{1}{3}$ of the validators violated one of the two Casper Commandments (a.k.a. slashing conditions).  These are:

\begin{enumerate}
   \item[\textbf{I.}] \textsc{A validator shalt not publish two or more nonidentical Prepares for same epoch.}
   
   In other words, a validator may prepare at most exactly one (\hash, \epochsource, \hashsource) triplet for any given epoch \epoch.

   \item[\textbf{II.}] \textsc{A validator shalt not publish a Commit between the epochs of a Prepare statement.} 
    
   Equivalently, a validator may not publish both $\langle \msgPREPARE, \epoch_p, \hash_p, \epochsource, \hashsource, \signature \rangle$ AND $\langle \msgCOMMIT, \epoch_c, \hash_c, \signature \rangle $ if $\epochsource < \epoch_c < \epoch_p$.

\end{enumerate}

If a validator violates a slashing condition, the evidence that they did this can be included into the blockchain as a transaction, at which point the validator's entire deposit will be taken away, with a 4\% ``finder's fee'' given to the submitter of the evidence transaction.

Finally, we define the ``ideal execution'' of the Casper protocol during an epoch $n$, as every validator preparing $C_{n}$ with $\hashsource = C_{n-1}$ and commiting $C_{n}$.  For example, during epoch $2$ (blocks $200 \ldots 299$), all validators prepare $b_{199}$ with $\hashsource = b_{99}$ and commit $b_{199}$.

\section{Proofs of Safety and Plausible Liveness}
\label{sect:theorems}

We prove Casper's two fundamental properties: \textit{accountable safety} and \textit{plausible liveness}. Accountable safety means that two conflicting checkpoints cannot be finalized unless $\geq \frac{1}{3}$ of validators violate a slashing condition (meaning at least one third of the total deposit is lost).  Plausible liveness means that, regardless of any previous events, it is always possible for $\frac{2}{3}$ of honest validators to finalize a new checkpoint.

\begin{figure}[h!tb]
\centering
    \includegraphics[width=5in]{conflicting_checkpoints.png}
\caption{Two checkpoints are \textit{conflicting} if they are on distinct chains, i.e. one is not an ancestor or a descendant of the other.}
\label{fig:conflicting_checkpoints}
\end{figure}


\begin{theorem}[Accountable Safety]
\label{theorem:safety}
Two conflicting checkpoints cannot be finalized unless $\geq \frac{1}{3}$ of validators violate a slashing condition.

\begin{proof}
Suppose the two conflicting checkpoints are $A$ in epoch $\epoch_A$ and $B$ in epoch $\epoch_B$. If both are finalized, this implies $\frac{2}{3}$ commits and $\frac{2}{3}$ prepares in epochs $\epoch_A$ and $e_B$. In the trivial case where $\epoch_A = \epoch_B$, this implies that some intersection of $\frac{1}{3}$ of validators must have violated slashing condition (1). In other cases, there must exist two chains $\Genesisblock < \ldots < \epoch_A^2 < \epoch_A^1 < \epoch_A$ and $\Genesisblock < \ldots < \epoch_B^2 < \epoch_B^1 < \epoch_B$ of justified checkpoints, both terminating at the genesis. Suppose without loss of generality that $\epoch_A > \epoch_B$. Then, there must be some $\epoch_A^i$ that either $\epoch_A^i = \epoch_B$ or $\epoch_A^i > \epoch_B > \epoch_A^{i+1}$. In the first case, since $A^i$ and $B$ both have $\frac{2}{3}$ prepares, at least $\frac{1}{3}$ of validators violated slashing condition (I). Otherwise, $B$ has $\frac{2}{3}$ commits and there exist $\frac{2}{3}$ prepares with $\epoch > B$ and $\epochsource < B$, so at least $\frac{1}{3}$ of validators violated slashing condition (II).
\end{proof}
\end{theorem}

\begin{figure}[h!tb]
\centering
   \begin{minipage}[b]{0.48\textwidth}
    \includegraphics[width=2.7in]{cs.pdf}
    \centering
    $\epoch_A = \epoch_B$
   \end{minipage}
   \begin{minipage}[b]{0.48\textwidth}
    \includegraphics[width=2.7in]{cs.pdf}
    \centering
    $\epoch_A < \epoch_B$
   \end{minipage}
\label{fig:conflicting_checkpoints}
\caption{Illustrating the two scenarios in Theorem \ref{theorem:safety}.}
\end{figure}

\begin{theorem}[Plausible Liveness]
\label{theorem:liveness}
It is always possible for $\frac{2}{3}$ of honest validators to finalize a new checkpoint, regardless of what previous events took place.

\begin{proof}
Suppose that all existing validators have sent some sequence of prepare and commit messages. Let $M$ with epoch $\epoch_M$ be the highest-epoch checkpoint that was justified, and let $n \ge \epoch_M$ be the highest epoch in which an honest validator prepared. Honest validators have not committed on any block which is not justified. Hence, neither slashing condition stops them from making prepares on a descendant of $M$ in epoch $n+1$, using $\epoch_M$ as $\epochsource$, and then committing this child.
\end{proof}

\end{theorem}

\section{Tweaking the Proposal Mechanism}
\label{sect:forkchoice}

Although Casper is chiefly an overlay on top of a proposal mechanism, in order to translate the \textit{plausible liveness} proven into the previous section into \textit{actual liveness in practice}, the proposal mechanism needs to be Casper-aware.  If the proposal mechanism isn't Casper-aware, as is the case with a proof of work chain that follows the typical fork-choice rule of ``always build atop the longest chain'', there are pathological scenarios where Casper gets ``stuck'' and any blocks built atop the longest chain cannot be finalized (or even justified) without violating a Commandment---one such pathological scenario is the fork in \figref{fig:forkchoice}.

In this case, $A_1$ or any descendant thereof cannot be finalized without slashing at least $\frac{1}{6}$ of validators. Finalizing $A_1$ itself would require $\frac{1}{3}$ of validators to violate Commandment \textbf{I} due to conflict with $B_1$, and finalizing any descendant would require $\frac{2}{3}$ prepares using $S$ as \hashsource, creating an intersection of size at least $\frac{1}{6}$ of validators that also committed $B_1$ and would thus violate Commandment \textbf{II}. However, miners on a proof of work chain would interpret $HASH1$ as the head and forever keep mining descendants of it, ignoring the chain based on $B_1$ which actually could get finalized.

\begin{figure}[h!tb]
\centering
\includegraphics[width=5.5in]{fork4.png}
\caption{Miners following the traditional proof of work fork choice rule would create blocks on HASH1, but because of the slashing conditions only blocks on top of HASH1' can be finalized.}
\label{fig:forkchoice}
\end{figure}

In fact, when \textit{any} checkpoint gets $k > \frac{1}{3}$ commits, no conflicting checkpoint can get finalized without $k - \frac{1}{3}$ of validators getting slashed. This necessitates modifying the proposal mechanism so that, in the event that some checkpoint is justified and receives any commits, any new blocks that are proposed are descendants of that checkpoint, and not conflicting checkpoints that have fewer or no commits.

In the case where the proposal mechanism is a proof of work chain, this entails modifying the fork choice rule: instead of blindly following a longest-chain rule, there needs to be an overriding rule that (i) finalized checkpoints are favored, and (ii) when there are no further finalized checkpoints, checkpoints with more (justified) commits are favored.

One complete description of such a rule would be:

\begin{enumerate}
\item Start with HEAD equal to the genesis of the chain.
\item Find the descendant checkpoint of HEAD with the most commits (only justified checkpoints are admissible). Set HEAD to this value.
\item Repeat (2) until no descendant with commits exists.
\item Choose the longest proof of work chain from there.
\end{enumerate}

The commit-following part of this rule can be viewed as mirroring the ``greegy heaviest observed subtree'' (GHOST) rule that has been proposed for proof of work chains\cite{sompolinsky2013accelerating}. \footnote{The symmetry is as follows. In GHOST, a node starts with the head at the genesis, then begins to move forward down the chain, and if it encounters a block with multiple children then it chooses the child that has the larger quantity of work built on top of it (including the child block itself and its descendants).

In this algorithm, we follow a similar approach, except we repeatedly seek the child closest to achieving finality. Commits on a descendant are implicitly commits on its ancestors, and so if a given descendant of a given block has the most commits, then we know that all children along the chain from the head to this descendant are closer to finality than any of their siblings; hence, looking for the \textit{descendant} with the most commits and not just the \textit{child} gives the right properties. Finalizing a checkpoint requires $\frac{2}{3}$ commits of a \textit{single} checkpoint, and so unlike GHOST we simply look for the maximum commit count instead of trying to sum up all commits in an entire subtree.}

This rule ensures that if there is a checkpoint such that no conflicting checkpoint can be finalized without at least some validators violating slashing conditions, then this is the checkpoint that will be viewed as the ``head'' and thus that validators will try to commit on. \footnote{Favoring checkpoint with even a single commit, instead of $\frac{1}{3} + \epsilon$, is desired to ensure that plausible liveness translates into actual liveness even in scenarios where up to $\frac{1}{3}$ of validators are offline.}

\section{Allowing Dynamic Validator Sets}
\label{sect:join_and_leave}

The set of validators needs to be able to change.  New validators must be able to join, and existing validators must be able to leave.  To accomplish this, we define a variable in the state called the \textit{dynasty counter}. When a would-be valdiator's deposit message is included in dynasty $d$, then the validator will be \textit{inducted} at the start of dynasty $d+2$.  We call $d+2$ this validator's \textit{start dynasty}.

The dynasty counter increments when the chain detects that the checkpoint of the current epoch that is part of its own history has been \textit{perfectly finalized} (that is, the checkpoint of epoch \epoch must be finalized during epoch \epoch, and the chain must learn about this before epoch \epoch ends). In simpler terms, when a user sends a ``deposit'' transaction, they need to wait for the transaction to be perfectly finalized, and then they need to wait again for the next epoch to be finalized; after this, they become part of the validator set.

For a validator to leave, they must send a ``withdraw'' message. If their withdraw message gets included during dynasty $n$, the validator similarly leaves the validator set during dynasty $n+2$; we call $n+2$ their \textit{end dynasty}. When a validator withdraws, their deposit is locked for a long period of time, the \textit{withdrawal delay} (think ``four months'') before they can take their money out. If, during the withdrawal delay, the validator violates any Commandment, the deposit is forfeited.

For a checkpoint to be justified, it must be prepared by a set of validators which contains (i) at least $\frac{2}{3}$ of the current dynasty (that is, validators with $startDynasty \le curDynasty < endDynasty$), and (ii) at least $\frac{2}{3}$ of the previous dyansty (that is, validators with $startDynasty \le curDynasty - 1 < endDynasty$. Finalization works similarly. The current and previous dynasties will usually greatly overlap; but in cases where they substantially diverge this ``stitching'' mechanism ensures that dynasty divergences do not lead to situations where a finality reversion or other failure can happen because different messages are signed by different validator sets and so equivocation is avoided.

\begin{figure}[h!tb]
\centering
\includegraphics[width=3.5in]{validator_set_misalignment.png}
\caption{Without the validator set stitching mechanism, it's possible for two conflicting checkpoints to be finalized with no validators slashed}
\label{fig:dynamic2}
\end{figure}

\subsection{Long Range Attacks}

Note that the withdrawal delay introduces a synchronicity assumption \textit{between validators and clients}. Because validators can withdraw their deposits after the withdrawal delay, there is an attack where a coalition of validators which had more than $\frac{2}{3}$ of deposits \textit{long ago in the past} withdraws their deposits, and then uses their historical deposits to finalize a new chain that conflicts with the original chain without fear of getting slashed. Despite violating slashing conditions to make a chain split, because the attacker has already withdrawn on both chains they do not lose any money. This is called the \textit{long-range atack}.

\begin{figure}
\centering
\includegraphics[width=3in]{LongRangeAttacks.png}
\caption{Despite violating slashing conditions to make a chain split, because the attacker has already withdrawn on both chains they do not lose any money. This is often called a \textit{long-range atack}.}
\label{fig:dynamic3}
\end{figure}

Suppose that clients can be relied on to log on at least once every time $\delta$ (think $\delta \approx$ 1 month). Then, if a client hears about two conflicting finalized checkpoints, $C_1$ at time $T_1$ and $C_2$ at time $T_2$, there are two cases. If, from the point of view of all clients, $T_2 > T_1$, then all clients will accept $C_2$ and there is no ambiguity. If on the other hand different clients see different orders (i.e. some see $T_2 > T_1$, others see $T_1 > T_2$), then it must be the case that for all clients, $|T_2 - T_1| \le 4\delta$. If the withdrawal delay is much greater than $4\delta$, then there is plenty of time during which slashing evidence can be included in both chains, and we make an assumption that if a chain does not include slashing evidence, then clients will reject this via a soft fork (see later section).

This rejection can also be partially automated, by implementing a mechanism where if a client sees a slashing condition violation before they see some block $B$ with timestamp $t$, and then the head is a descendant of $B$ with timestamp greater than $t + EIT$ (``evidence inclusion timelimit'') that does has not yet punished the malfeasant validator, the client shows a warning.

\section{Recovering from Castastrophic Crashes}
\label{sect:leak}

Suppose that $>\frac{1}{3}$ of validators crash-fail at the same time---i.e, they are no longer connected to the network due to a network partition, computer failure, or are malicious actors. Then, no later checkpoint will be able to get finalized.

We can recover from this by instituting a ``leak'' which dissipates the deposits of validators that do not prepare or commit, until eventually their deposit sizes decrease low enough that the validators that \textit{are} preparing and committing become a $\frac{2}{3}$ supermajority. The simplest possible formula is something like ``validators with deposit size $D$ lose $D * p$ in every epoch in which they do not prepare and commit'', though to resolve catastrophic crashes more quickly a formula which increases the rate of dissipation in the event of a long streak of non-finalized blocks may be optimal.

The dissipated portion of deposits can either be burned or simply forcibly withdrawn and immediately refunded to the validator; which of the two strategies to use, or what combination, is an economic incentive concern and thus outside the scope of this paper.

Note that this does introduce the possibility of two conflicting checkpoints being finalized, with validators only losing money on one of the two checkpoints as seen in \figref{fig:commitsync}.

\begin{figure}[h!tb]
\centering
\includegraphics[width=4in]{CommitsSync.png}
\caption{The checkpoint on the left can be finalized immediately. The checkpoint on the right can be finalized after some time, once offline validator deposits sufficiently dissipate.}
\label{fig:commitsync}
\end{figure}

If the goal is simply to achieve maximally close to 50\% fault tolerance, then clients should simply favor the finalized checkpoint that they received earlier. However, if clients are also interested in defeating 51\% censorship attacks, then they may want to at least sometimes choose the minority chain. All forms of ``51\% attacks'' can thus be resolved fairly cleanly via ``user-activated soft forks'' that reject what would normally be the dominant chain. Particularly, note that finalizing even one block on the dominant chain precludes the attacking validators from preparing on the minority chain because of Commandment II, at least until their balances decrease to the point where the minority can commit, so such a fork would also serve the function of costing the majority attacker a very large portion of their deposits.

\section{Conclusions}

This introduces the basic workings of Casper the Friendly Finality Gadget's prepare and commit mechanism and fork choice rule, in the context of Byzantine fault tolerance analysis. Separate papers will serve the role of explaining and analyzing incentives inside of Casper, and the different ways that they can be parametrized and the consequences of these paramtrizations.

\section{Acknowledgements}

We thank Virgil Griffith for review and Sandro Lera for mathematics.


%\section{References}
\bibliographystyle{abbrv}
\bibliography{ethereum}

%\appendix
\clearpage
\section{Unused text}

\todo{This is where text goes that for which a home hasn't been found yet.  If no home is found, it will be deleted.}



We define an \textit{epoch} as a range of 100 blocks (e.g., blocks 600...699 are epoch 6), and a \textit{checkpoint} of an epoch is the final block in that epoch. 

The proposal mechanism will initially be the existing Ethereum proof of work chain, making the first version of Casper a \textit{hybrid PoW/PoS algorithm} that relies on proof of work for liveness but not safety, but in future versions the proposal mechanism can be substituted with something else.



for the same \epoch and \hash as in \eqref{eq:msgPREPARE}.  The $\hash$ is the block hash of the block at the start of the epoch.  A hash $\hash$ being justified entails that all fresh (non-finalized) ancestor blocks are also justified.  A hash $\hash$ being finalized entails that all ancestor blocks are also finalized, regardless of whether they were previously fresh or justified.  An ``ideal execution'' of the protocol is one where, at the start of every epoch, every validator Prepares and Commits the first blockhash of each epoch, specifying the same $\epochsource$ and $\hashsource$.

In the Casper protocol, there exists a set of validators, and in each \textit{epoch} (see below) validators may send two kinds of messages: $$[PREPARE, epoch, hash, epoch_{source}, hash_{source}]$$ and $$[COMMIT, epoch, hash]$$


If, during an epoch $e$, for some specific ancestry hash $h$, for any specific ($epoch_{source}, hash_{source}$ pair), there exist $\frac{2}{3}$ prepares of the form $$[PREPARE, e, h, epoch_{source}, hash_{source}]$$, then $h$ is considered \textit{justified}. If $\frac{2}{3}$ commits are sent of the form $$[COMMIT, e, h]$$ then $h$ is considered \textit{finalized}.


During epoch $n$, validators are expected to send prepare and commit messages with $\epoch = n$ and $h$ equal to a checkpoint of epoch $n$. Prepare messages may specify as \hashsource a checkpoint for any previous epoch (preferably the preceding checkpoint) of \hash, and which is \textit{justified} (see below), and the \epochsource is expected to be the epoch of that checkpoint.


Honest validators will never violate slashing conditions, so this implies the usual Byzantine fault tolerance safety property, but expressing this in terms of slashing conditions means that we are actually proving a stronger claim: if two conflicting checkpoints get finalized, then at least $\frac{1}{3}$ of validators were malicious, \textit{and we know whom to blame, and so we can maximally penalize them in order to make such faults expensive}.


This simplifies our finalty mechanism because it allows it to be expressed as a fork choice rule where the ``score'' of a block only depends on the block and its children, putting it into a similar category as more traditional PoW-based fork choice rules such as the longest chain rule and GHOST\cite{sompolinsky2013accelerating}.

Unlike GHOST, however, this fork choice rule is also \textit{finality-bearing}: there exists a ``finality'' mechanism that has the property that (i) the fork choice rule always prefers Finalized blocks over non-Finalized blocks, and

\begin{figure}[h!tb]
\centering
    \includegraphics[width=3in]{prepares_commits.png}
%    \includegraphics[width=5.5in]{cs.pdf}
    \caption{Illustrating prepares, commits and checkpoints. Arrows represent \textit{dependency} (e.g., a commit depends on there being $\frac{2}{3}$ existing prepares).}
    \label{fig:prepares_and_commits}
\end{figure}


\begin{enumerate}
\item Start with HEAD equal to the genesis of the chain.
\item Select the descendant checkpoint of HEAD with the most commits (only Justified checkpoints are admissible)
\item Repeat (2) until no descendant with commits exists.
\item Choose the longest proof of work chain from there.
\end{enumerate}



\section{Notes to Authors}
\subsection{Questions}
\begin{itemize}
\item True/False: The Dynasty counter increments iff there's been a finalization?
\end{itemize}


\subsection{Notes on Suggested Terminology}
\begin{itemize}
\item parent $\rightarrow$ predecessor.
\item child $\rightarrow$ successor (unless want to emphasize there can be multiple candidate successors)
\item ancestors $\rightarrow$ lineage
\item to refer to the set of $\{$ predecessor, successor $\}$ $\rightarrow$ adjacent
\end{itemize}


\subsection{Todo}
\begin{itemize}
\item \todo{Reference the various Figures within the text so we more easily know what goes with what.}
\item \todo{fill me in}
\end{itemize}


In the other way...



%%%%%%%%%%%%%%%%%%%%%%%%%%%%%%%%%%%%%%%%%%%%%%%%%%%%%%


\end{document}
        
